\title{QLoRA: Efficient Finetuning of Quantized LLMs}

\begin{abstract}
We introduce \textsc{QLoRA}, a memory-efficient finetuning method that enables a 65B parameter language model to be finetuned on a single 48GB GPU, while achieving performance comparable to standard 16-bit finetuning. \textsc{QLoRA} enables gradient backpropagation through a frozen, 4-bit quantized version of a pretrained language model into Low Rank Adapters~(LoRA). The resulting models, which we refer to as \textbf{Guanaco}, surpass all previously released open-source models on the Vicuna benchmark, attaining 99.3\% of ChatGPT's performance after just 24 hours of single-GPU finetuning. The efficiency of \textsc{QLoRA} is driven by several novel memory-saving techniques: (a) 4-bit NormalFloat (NF4), a newly proposed datatype that is theoretically optimal for weights with a normal distribution; (b) Double Quantization, which further lowers memory consumption by quantizing the quantization parameters themselves; and (c) Paged Optimizers, which efficiently handle memory spikes. Utilizing \textsc{QLoRA}, we finetune over 1,000 models and conduct comprehensive evaluations of instruction following and chatbot abilities across 8 instruction datasets, diverse model architectures (including LLaMA and T5), and scales not tractable for conventional finetuning (such as 33B and 65B models). The findings demonstrate that finetuning with QLoRA on a compact, high-quality dataset achieves state-of-the-art outcomes, even for smaller models compared to the former SoTA. We provide in-depth analyses based on both human judgment and GPT-4 automated evaluations, and show that GPT-4-based assessments offer a practical and cost-effective alternative to human annotation. Moreover, our work indicates significant limitations in the reliability of existing chatbot benchmarks. A detailed, failure-focused review highlights specific shortcomings of \textbf{Guanaco} relative to ChatGPT. We openly share our models and codebase, which includes CUDA kernels supporting 4-bit training.\footnote{\url{https://github.com/artidoro/qlora} and \url{https://github.com/TimDettmers/bitsandbytes}}
\end{abstract}

\section{Introduction}

Adapting large language models (LLMs) through finetuning has proven to be a robust strategy for enhancing their capabilities \citep{min2021metaicl, wei2021finetuned, ouyang2022training, wang2022super, wang2022self, liu2022few}, as well as for incorporating or eliminating specific behaviors as desired \citep{ouyang2022training,askell2021general,bai2022training}. Nevertheless, performing finetuning on extremely large models remains prohibitively costly: for example, standard 16-bit finetuning of the LLaMA model with 65 billion parameters~\citep{touvron2023llama} consumes upwards of 780 GB of GPU memory. Although state-of-the-art quantization approaches can significantly lower LLM memory usage during inference \citep{dettmers2022llm, dettmers2022case, frantar2022gptq, xiao2022smoothquant}, they generally fail to extend these memory savings to the training process \citep{wortsman2023stable}.

Here, we introduce a novel approach demonstrating, for the first time, that finetuning a 4-bit quantized model is feasible without incurring any drop in performance. Our technique, \textsc{QLoRA}, first applies an innovative high-precision quantization method to compress a pretrained model to 4 bits and subsequently incorporates a compact set of trainable Low-rank Adapter weights \citep{hu2021lora} 

that can be updated via backpropagation through the quantized parameters.

With \textsc{QLoRA}, the average memory consumption required for finetuning a 65B parameter model decreases drastically, from over 780GB to less than 48GB of GPU memory, all while retaining equivalent runtime efficiency and predictive power relative to the standard 16-bit finetuning baseline. This advancement dramatically increases the accessibility of LLM finetuning, enabling the finetuning of the largest open models on a single GPU. Employing \textsc{QLoRA}, we develop the \textbf{Guanaco} model series, where the second largest model achieves 97.8\% of ChatGPT's performance on the Vicuna~\citep{vicuna2023} evaluation, training in under 12 hours with a single consumer GPU. When leveraging a single professional GPU for 24 hours, our largest configuration attains 99.3\%, effectively matching ChatGPT’s performance on Vicuna. Notably, during inference, the smallest \textbf{Guanaco} variant (7B parameters) operates with only 5 GB of memory and surpasses a 26 GB Alpaca model by more than 20 percentage points on the Vicuna benchmark (Table~\ref{tbl:vicuna_eval}).

\textsc{QLoRA} incorporates several key advancements aimed at minimizing memory consumption without compromising performance: (1) \textbf{4-bit NormalFloat}, an information-theoretically optimal quantization scheme tailored for normally distributed values, which empirically surpasses both 4-bit Integer and 4-bit Float formats. (2) \textbf{Double Quantization}, which further compresses memory usage by quantizing the quantization constants themselves, resulting in a mean savings of about 0.37 bits per parameter—translating to around 3 GB for models with 65B parameters. (3) \textbf{Paged Optimizers}, which leverage NVIDIA’s unified memory to mitigate the memory surges linked to gradient checkpointing during training on mini-batches with extended sequence lengths. These elements are integrated to yield a refined LoRA approach, embedding adapters in all network layers, thus substantially reducing the typical accuracy tradeoffs observed in previous techniques.

By increasing efficiency, \textsc{QLoRA} allows for comprehensive empirical investigation of instruction finetuning and chatbot capabilities at scales infeasible for standard finetuning approaches due to memory constraints. Consequently, we fine-tune over 1,000 distinct models, spanning diverse instruction tuning datasets, model types, and ranging from 80M up to 65B parameters. Our results indicate that \textsc{QLoRA} restores performance to the level of 16-bit finetuning (\S\ref{sec:comp_to_std}), and facilitates the development of a leading chatbot, \textbf{Guanaco}, (\S\ref{sec:pushing}). Furthermore, our analysis reveals two important patterns in the trained models. First, dataset quality considerably outweighs mere dataset size; for instance, a small dataset of 9,000 examples (OASST1) outperforms a much larger dataset containing 450,000 examples (FLAN v2, subsampled) on chatbot evaluation metrics, even when both target generalization to instruction-following. Second, we demonstrate that excelling in the Massive Multitask Language Understanding (MMLU) benchmark does not necessarily entail strong results on the Vicuna chatbot benchmark, and vice versa; this underscores that the appropriateness of the dataset for the specific task is more critical than its scale.

Additionally, we conduct a comprehensive examination of chatbot performance, incorporating assessments from both human annotators and GPT-4. Our evaluation employs a tournament-style framework where models are pitted against one another, each aiming to generate the most appropriate response to a given prompt. In each round, either GPT-4 or human raters determine the superior response. Results from these competitions are synthesized into Elo scores~\citep{elo1967proposed, elo1978rating}, enabling us to objectively rank the chatbots. Our analysis indicates a high level of concordance between GPT-4-based and human judgment in ranking model performance; however, notable exceptions exist where the two methods diverge significantly. Consequently, we emphasize that, although model-driven evaluation offers a more economical alternative to human annotation, it also introduces additional sources of uncertainty.

To complement the quantitative results of our chatbot benchmarking, we present an in-depth qualitative study of the \textbf{Guanaco} models. This investigation brings to light exemplary successes as well as failure modes that are not reflected in aggregate benchmark scores.

In support of transparency and reproducibility, we publicly release all model outputs, along with the accompanying human and GPT-4 annotations. Furthermore, we make our entire codebase and CUDA kernels available as open source, and have integrated our methods into the Hugging Face transformers framework \citep{wolf2019huggingface} for straightforward access. We provide a set of adapter modules for models of size 7/13/33/65B, each trained on one of eight instruction-following datasets, culminating in a total of 32 distinct finetuned, open-source models.

\section{Background}
\paragraph{Block-wise k-bit Quantization} Quantization refers to the transformation of input data from a format with high representational fidelity to one with lower fidelity. Typically, this involves converting higher-bitwidth numerical types, such as 32-bit floating point numbers, into types with fewer bits, such as 8-bit integers. To maximize utilization of the reduced data type’s dynamic range, input tensors are often normalized by their maximum absolute value before conversion. As an example, converting a 32-bit floating point (FP32) tensor into an Int8 tensor within $[-127, 127]$ proceeds as follows:

\begin{equation}
    \mathbf{X}^{ \text{Int8}} = \text{round}\left( \frac{127}{{\text{absmax}}(\mathbf{X}^{{\text{FP32}}})}\mathbf{X}^{{\text{FP32}}} \right) = \text{round}(c^{\text{FP32}}\cdot \mathbf{X}^{{\text{FP32}}}),
\end{equation}

where $c$ is called the {\em quantization constant} or {\em quantization scale}. The reverse transformation, dequantization, is performed via:

\begin{equation}
    \text{dequant}(c^{\text{FP32}}, \mathbf{X}^{\text{Int8}}) = \frac{\mathbf{X}^{\text{Int8}}}{c^{\text{FP32}}} = \mathbf{X}^{\text{FP32}}
\end{equation}

A limitation of this method arises when an outlier—an unusually large value—appears in the input tensor, causing the quantization bins (specific combinations of bits) to be inefficiently populated; some bins may have very few or no values mapped to them. To address this issue, it is standard practice to partition the input tensor into blocks, each quantized separately with its distinct quantization constant $c$. More specifically, the input tensor $\mathbf{X}\in \mathbb{R}^{b\times h}$ is divided into $n$ consecutive segments of length $B$ by flattening the tensor and segmenting it linearly, resulting in ${n = ({b\times h} )/{B} }$ individual blocks. Each block is then quantized independently, following Equation~1, producing a quantized version of the tensor along with $n$ separate quantization constants $c_i$.

\paragraph{Low-rank Adapters} The LoRA (Low-rank Adapter) approach to finetuning \citep{hu2021lora} reduces memory overhead by introducing a limited number of trainable parameters, called adapters, and retaining the bulk of the model parameters in their pretrained state. During training, stochastic gradient descent backpropagates gradients through the frozen model weights to these adapters, which are updated to minimize the loss. LoRA modifies a linear transformation by supplementing it with a factorized, low-rank projection. Considering a linear transformation ${\mathbf{X} \mathbf{W} = \mathbf{Y}}$, where $\mathbf{X}\in  \mathbb{R}^{b\times h}$ and $\mathbf{W}\in  \mathbb{R}^{h\times o}$, the LoRA method computes:

\begin{equation}
    \mathbf{Y} = \mathbf{X}\mathbf{W} + s\mathbf{X}\mathbf{L}_1\mathbf{L}_2,
\end{equation}

where the matrices $\mathbf{L}_1\in  \mathbb{R}^{h\times r}$ and $\mathbf{L}_2\in  \mathbb{R}^{r\times o}$, and $s$ denotes a scalar multiplier.

\paragraph{Memory Requirement of Parameter-Efficient Finetuning} A key consideration concerns the memory usage associated with LoRA during training, particularly in relation to both the count and size of the adapters employed. Due to LoRA's notably small memory demands, deploying additional adapters to boost model performance is feasible without substantially increasing total memory consumption. Although LoRA is categorized as a Parameter Efficient Finetuning (PEFT) strategy, the dominant contributor to memory usage in large language model (LLM) finetuning is typically the storage of activation gradients, not the learned parameters introduced by LoRA. For instance, in the context of training a 7B LLaMA model on FLAN v2 with a batch size of 1, where LoRA weights account for roughly 0.2\% of the base model parameters\citep{hu2021lora,liu2022few}, input gradients from LoRA utilize 567 MB of memory, while the LoRA parameter tensors themselves occupy just 26 MB. Implementing gradient checkpointing~\citep{chen2016training} lowers the average per-sequence input gradient memory usage to 18 MB, yet this figure still surpasses the combined size of all LoRA weights. In stark contrast, the underlying 4-bit model demands 5,048 MB. These findings underscore the value of gradient checkpointing and indicate that further minimizing the LoRA parameter size offers only limited reductions in memory load. Consequently, adding extra adapters remains viable without a dramatic impact on overall training memory requirements (for further details, refer to Appendix~\ref{app:memory}). As will be elaborated later, this capacity is essential to achieving performance that matches full 16-bit precision.

\section{\textsc{QLoRA} Finetuning}

\textsc{QLoRA} facilitates effective finetuning at 4-bit precision through the introduction of two principal techniques: 4-bit NormalFloat (NF4) quantization and Double Quantization. The method further incorporates Paged Optimizers, a mechanism designed to mitigate out-of-memory problems frequently caused by memory surges during gradient checkpointing—thus addressing a major obstacle in large-scale model finetuning on a single machine.

The \textsc{QLoRA} framework utilizes a low-precision storage data type—typically 4-bit in our experiments—and a separate data type for computations, commonly BFloat16. Operationally, this involves dequantizing \textsc{QLoRA} weight tensors to BFloat16 whenever computations are needed, after which the associated matrix multiplications are performed at 16-bit precision.

We next detail the constituent mechanisms within \textsc{QLoRA}, leading to a precise mathematical description of the method.

\paragraph{4-bit NormalFloat Quantization} The NormalFloat (NF) data type extends Quantile Quantization \citep{dettmers2022optimizers}, a method that is theoretically optimal for information preservation and guarantees that quantization bins each capture the same number of tensor elements from the input. Quantile quantization is realized by computing the quantiles of the input tensor using its empirical cumulative distribution function.

However, a notable drawback is the significant computational cost associated with quantile estimation. To alleviate this, rapid approximation strategies—such as SRAM quantiles~\citep{dettmers2022optimizers}—are employed. Nevertheless, the inherent approximations introduce substantial quantization errors, especially for outliers, which tend to carry high informational value.

When input tensors derive from a distribution invariant under quantization scaling, the expense and errors associated with quantile estimation can be bypassed. In these scenarios, shared quantiles across tensors permit efficient and exact quantile calculation.

Pretrained neural network weights are typically distributed according to a normal distribution centered at zero with standard deviation $\sigma$ (refer to Appendix~\ref{app:normality}). To standardize all weights into a consistent reference distribution, we scale $\sigma$ such that the resulting distribution aligns perfectly with the interval specified by our chosen data type. For this work, we adopt $[-1, 1]$ as our data type’s range. This requires that both the quantization levels of the data type and the weight values themselves are adjusted to fall within $[-1, 1]$.

For normal distributions with mean zero and any standard deviation $\sigma$ within $[-1, 1]$, the data type achieving optimal information-theoretic efficiency is determined by the following procedure: (1) compute $2^k+1$ quantile positions of an idealized $N(0, 1)$ distribution, thereby constructing a $k$-bit quantile-based quantizer tailored to normal distributions; (2) map these quantization points into the $[-1, 1]$ interval by normalization; and (3) transform the input weight tensor into the $[-1, 1]$ range using rescaling based on its absolute maximum value.

With the scales of the weights and the data type matched, standard quantization techniques can then be applied. Notably, the normalization performed in step (3) serves to adjust the standard deviation of the weight tensor so it aligns with that of the k-bit data type. More formally, the $2^k$ quantization values $q_i$ for the data type can be computed as:

\begin{equation}
q_i  = \frac{1}{2}\left( Q_X\left(\frac{i}{2^k+1}\right) + Q_X\left(\frac{i+1}{2^k+1}\right)\right),
\end{equation}

where $Q_X(\cdot)$ denotes the quantile function of the standard normal distribution $N(0,1)$. A notable limitation of symmetric k-bit quantization is its inability to represent zero exactly, which is essential for losslessly quantizing zero-valued entities such as padding. To guarantee that zero maps precisely to a quantization bin and to fully utilize all $2^k$ quantization levels, we devise an asymmetric quantization scheme. Specifically, we compute quantiles $q_i$ for two intervals: $2^{k-1}$ quantiles for negative values and $2^{k-1}+1$ for non-negative values, subsequently merging these sets and eliminating the duplicate zero entry. We refer to the resulting quantization format as \emph{k-bit NormalFloat} (NFk), which distributes data so that each quantization bin contains an equal expected frequency and provides information-theoretic optimality for zero-centered Gaussian distributions. Details of the precise bin values are included in Appendix~\ref{app:nf4}.

\paragraph{Double Quantization{}} 
We propose \emph{Double Quantization} (DQ), wherein quantization constants themselves are quantized for further reductions in memory usage. While accurate 4-bit quantization demands small blocksizes~\citep{dettmers2022case}, this design choice introduces non-negligible memory cost for storing quantization constants. For instance, when 32-bit constants and a blocksize of 64 are employed to quantize matrix $\mathbf{W}$, these constants incur an average overhead of $32/64 = 0.5$ bits per model parameter. Double Quantization aims to lower this memory requirement.

Concretely, Double Quantization performs a secondary quantization on the set of quantization constants $c_2^{\text{FP32}}$ produced by the primary quantization, generating a more compressed representation $c_2^{\text{FP8}}$ and additional quantization constants $c_1^{\text{FP32}}$ for this second stage. For the secondary quantization, we apply 8-bit Floating Point representation with a blocksize of 256, leveraging empirical findings from \citet{dettmers2022case} that confirm no observed degradation for 8-bit quantization. Because $c_2^{\text{FP32}}$ is strictly positive, we recenter its values by subtracting the mean, allowing the use of symmetric quantization. This approach lowers the mean memory required per parameter for a blocksize of 64 from $0.5$ bits ($32/64$) down to $0.127$ bits ($8/64 + 32/(64\cdot256)$), resulting in an average savings of 0.373 bits per parameter.

\paragraph{Paged Optimizers} leverage NVIDIA's unified memory~\footnote{\tiny \url{https://docs.nvidia.com/cuda/cuda-c-programming-guide}} capability, which transparently manages memory pages between the CPU and GPU. This automatic handling allows seamless execution even when GPU memory is temporarily exhausted, mirroring the mechanism of conventional paging between main memory and storage. We exploit this unified memory to assign optimizer state buffers in paged memory, which are then offloaded to CPU RAM when GPU memory becomes scarce and subsequently brought back to GPU as needed during optimizer updates.

\paragraph{\textsc{QLoRA}.} Incorporating the aforementioned mechanisms, we specify \textsc{QLoRA} for a quantized single linear layer equipped with one LoRA adapter as:

\begin{equation}
    \mathbf{Y}^{\text{BF16}} =  \mathbf{X}^{\text{BF16}} \text{doubleDequant}(c_1^{\text{FP32}}, c_2^{\text{k-bit}}, \mathbf{W}^{\text{NF4}}) + \mathbf{X}^{\text{BF16}}\mathbf{L}^{\text{BF16}}_1\mathbf{L}^{\text{BF16}}_2,
\end{equation}

where doubleDequant$(\cdot)$ is formulated as:

\begin{equation}
    \text{doubleDequant}(c_1^{\text{FP32}}, c_2^{\text{k-bit}}, \mathbf{W}^{\text{k-bit}}) = \text{dequant}(\text{dequant}(c_1^{\text{FP32}}, c_2^{\text{k-bit}}), \mathbf{W}^{\text{4bit}}) = \mathbf{W}^{\text{BF16}},
\end{equation}

Here, $\mathbf{W}$ utilizes the NF4 format, while $c_2$ is stored in FP8.

To achieve greater quantization fidelity, a blocksize of 64 is chosen for $\mathbf{W}$. For $c_2$, a blocksize of 256 is used to reduce memory demands.

During parameter optimization, it suffices to compute the gradients with respect to the adapter weights $\frac{\partial E}{\partial \mathbf{L}_i}$, with no requirement to update the gradients for the 4-bit weights $\frac{\partial E}{\partial \mathbf{W}}$. Nevertheless, evaluating $\frac{\partial E}{\partial \mathbf{L}_i}$ implicitly requires calculating $\frac{\partial \mathbf{X}}{\partial \mathbf{W}}$, which is accomplished through equation~(5), involving the dequantization process from the storage format $\mathbf{W}^{\text{NF4}}$ to the computation data type $\mathbf{W}^{\text{BF16}}$ in order to compute the derivative $\frac{\partial \mathbf{X}}{\partial \mathbf{W}}$ at BFloat16 precision.

In summary, \textsc{QLoRA} employs a single data type for storage (typically 4-bit NormalFloat) and a separate data type for computation (16-bit BrainFloat). During both the forward and backward passes, the storage data is dequantized to the computation data type; however, gradient updates are computed exclusively for the LoRA parameters, which are maintained in 16-bit BrainFloat precision.

\section{QLoRA vs. Standard Finetuning}
\label{sec:comp_to_std}

Having outlined the mechanics of QLoRA and its ability to considerably decrease memory demands during model finetuning, the crucial question is whether QLoRA can achieve comparable results to full-model finetuning. Moreover, we seek to dissect various components of QLoRA, including the effect of utilizing NormalFloat4 as opposed to conventional Float4. The subsequent sections detail experiments conducted to address these issues.

\paragraph{Experimental setup.} Three model types—encoder, encoder-decoder, and decoder-only—are considered for evaluating QLoRA in comparison to 16-bit adapter-finetuning and full-model finetuning, focusing on model sizes up to 3B. Our assessment framework involves the GLUE benchmark \citep{wang2018glue} with RoBERTa-large \citep{liu2019roberta}, Super-NaturalInstructions (TKInstruct) \citep{wang2022super} with T5 \citep{t5}, as well as 5-shot MMLU \citep{hendrycksmeasuring} using LLaMA finetuned on Flan v2 \citep{longpre2023flan} and Alpaca \citep{alpaca}. To further investigate the merits of NF4 compared to alternative 4-bit quantization formats, we employ the evaluation strategy described by \citet{dettmers2022case}, measuring post-quantization zero-shot accuracy and perplexity on a range of models (OPT~\citep{zhang2022opt}, LLaMA~\citep{touvron2023llama}, BLOOM~\citep{scao2022bloom}, Pythia~\citep{biderman2023pythia}) covering parameter counts from 125m to 13B.

Additional details for each experiment are included in the results section to enhance clarity. Comprehensive experimental information is presented in Appendix~\ref{app:exp-setup}.

Paged optimizers are indispensable for conducting \textsc{QLoRA} finetuning of 33B/65B models on single GPUs with 24GB or 48GB memory; however, we do not report explicit quantitative metrics for Paged Optimizers, as paging occurs primarily with long-sequence mini-batches, which are infrequent. Nevertheless, we analyze the training runtime of paged optimizers with 65B models on 48GB GPUs and observe that, at a batch size of 16, their performance is on par with standard optimizers. Future investigations should systematically evaluate and characterize potential slow-downs arising from paging.

\paragraph{Default LoRA hyperparameters do not match 16-bit performance}

Applying the conventional approach of targeting LoRA to the query and value attention projection matrices \citep{hu2021lora} does not recover full-model finetuning performance on larger base models. As evidenced by Figure~\ref{fig:hyper} for Alpaca finetuning on LLaMA 7B, our findings indicate that the total number of LoRA adapters is the most influential hyperparameter, and only when LoRA is applied across all linear transformer block layers does performance approach that of full finetuning. Other parameters, such as the projection rank $r$, exert minimal impact on outcomes (see Appendix \ref{app:exp-setup}).

In a similar vein, our analysis reveals that the standard hyperparameter settings used for fully finetuned baseline models are insufficiently optimized. To address this, we perform an extensive grid search across learning rates ranging from $1 \times 10^{-6}$ to $5 \times 10^{-5}$, and batch sizes from 8 to 128, in order to establish robust baseline results. The outcomes for 7B LLaMA finetuned on the Alpaca dataset are illustrated in Figure~\ref{fig:hyper}.

\paragraph{4-bit NormalFloat surpasses 4-bit Floating Point in empirical performance}

Although the 4-bit NormalFloat (NF4) format is theoretically optimal from an information perspective, its practical benefits remain to be substantiated. Adopting the experimental methodology from \citet{dettmers2022case}, we evaluate various quantized LLMs (OPT \citep{zhang2022opt}, BLOOM \cite{scao2022bloom}, Pythia \cite{biderman2023pythia}, LLaMA), spanning sizes from 125M to 65B parameters, using different data representations on both language modeling benchmarks and several zero-shot tasks. As shown in Figure~\ref{fig:nf4} and Table~\ref{tbl:nf4_ppl}, NF4 consistently demonstrates a marked improvement over FP4 and Int4, and employing double quantization achieves significant reductions in memory usage without negatively impacting model accuracy.

\paragraph{k-bit \textsc{QLoRA} achieves parity with 16-bit full finetuning and 16-bit LoRA approaches} 

Prior research has indicated that although 4-bit quantization is viable for \emph{inference}, it often results in diminished performance compared to 16-bit representations \citep{dettmers2022case,frantar2022gptq}. This prompts an important inquiry: can the apparent drop in performance be offset by subsequently applying 4-bit adapter finetuning? To investigate this, we design two experimental conditions.

The first setting involves comparing the results of full 16-bit finetuning against those of adapter-based finetuning at 16-bit, 8-bit, and 4-bit precision, using RoBERTA and T5 models with sizes ranging from 125M to 3B parameters. The evaluation takes place on both GLUE and the Super-NaturalInstructions datasets, as reported in Table~\ref{tab:nf4vsbf16}. In both cases, adapter-based methods at varying precisions effectively match the performance of the 16-bit fully finetuned baseline, suggesting that any accuracy lost due to quantization can be completely regained through appropriate adapter finetuning post-quantization.

In our second experimental configuration, due to the requirement for multiple high-memory GPU servers to fully finetune models with 11B parameters or more, we investigate whether 4-bit \textsc{QLoRA} achieves parity with 16-bit LoRA across models ranging from 7B to 65B parameters. Specifically, we conduct finetuning on LLaMA models of 7B through 65B using the Alpaca and FLAN v2 instruction following datasets, subsequently evaluating these models with the MMLU benchmark employing a 5-shot accuracy metric. Table~\ref{tbl:llama-datatype} presents these results, illustrating that the NF4 configuration with double quantization can entirely replicate the MMLU performance achieved by 16-bit LoRA. Furthermore, we observe that \textsc{QLoRA} with FP4 demonstrates approximately a 1 percentage point lower performance relative to the 16-bit brain float LoRA baseline. These findings substantiate our earlier conclusions: (1) \textsc{QLoRA} equipped with NF4 attains results equivalent to both 16-bit full finetuning and 16-bit LoRA, and (2) NF4 delivers higher quantization fidelity than FP4.

\paragraph{Summary} The collective results from our experiments robustly establish that 4-bit \textsc{QLoRA} using the NF4 data format delivers comparable performance to both 16-bit full finetuning and 16-bit LoRA on academic evaluation benchmarks. Additionally, our analyses confirm that NF4 surpasses FP4 in effectiveness, and that utilizing double quantization does not impair finetuning outcomes. Taken together, these insights provide strong evidence supporting the reliability of 4-bit \textsc{QLoRA} in achieving results on par with standard 16-bit finetuning techniques.

Consistent with trends observed in previous quantization research \citep{dettmers2022case}, our experimental findings on MMLU and Elo suggest that, under fixed finetuning and inference resource constraints, it is advantageous to expand the size of the base model while utilizing lower numerical precision. This observation underscores the efficiency advantages rendered by \textsc{QLoRA}. Notably, our experiments reveal no drop in performance with 4-bit finetuning compared to full-finetuning, raising a pertinent open question about the precise boundary of the performance-precision trade-off in QLoRA, a topic we propose for future investigation.

Our investigation turns to the feasibility of instruction tuning at model and dataset scales that exceed the capacity of traditional 16-bit finetuning on standard academic computing resources.

\section{Pushing the Chatbot State-of-the-art with QLoRA}
\label{sec:pushing}
After demonstrating that 4-bit \textsc{QLoRA} achieves performance comparable to conventional 16-bit finetuning across a variety of scales, tasks, and datasets, we systematically examine instruction tuning on the largest open-source language models currently accessible to researchers. To gauge the effectiveness of these models under instruction tuning, we utilize a demanding Natural Language Understanding benchmark (MMLU), and further introduce novel techniques for evaluating chatbot behavior in practical, real-world contexts.

\subsection{Experimental setup}
An outline of our experimental framework is provided here, with comprehensive specifics available in Appendix \ref{app:chatbot}.

\paragraph{Data} In light of the absence of an extensive review of contemporary instruction-following datasets, we select eight recently developed datasets for this work. These datasets encompass various collection strategies: crowd-sourced data (OASST1~\citep{kopf2023openassistant}, HH-RLHF~\citep{bai2022training}), distillations from models previously fine-tuned on instructions (Alpaca~\citep{alpaca}, self-instruct~\citep{wang2022self}, unnatural-instructions~\citep{honovich2022unnatural}), aggregated corpora (FLAN v2~\citep{chung2022scaling}), as well as datasets blending multiple approaches (Chip2~\citep{laion2023}, Longform~\citep{koksal2023longform}). Collectively, these datasets offer a spectrum of languages, sizes, and licensing arrangements.

\paragraph{Training Setup} To eliminate possible confounding variables from heterogeneous training methodologies, we carry out QLoRA instruction finetuning exclusively with a cross-entropy loss in a supervised learning setup, intentionally omitting reinforcement learning approaches, even when the dataset contains human evaluation signals. For resources that clearly delineate instructions and corresponding responses, finetuning is limited to the response segment only (refer to ablation studies in Appendix \ref{app:chatbot}). For OASST1 and HH-RLHF, where multiple candidate responses exist, the highest-rated response from each dialogue tree node is selected, and finetuning is performed over the entire chosen conversation, inclusive of instructions. Every experiment utilizes NF4 \textsc{QLoRA} equipped with double quantization and paged optimizer mechanisms to avert memory usage surges during gradient checkpointing. We undertake modest hyperparameter tuning on LLaMA 13B and 33B; all parameters optimized for 7B models were transferrable—apart from the learning rate and batch size. For 33B and 65B models, we set the learning rate to half its original value and double the batch size.

\paragraph{Baselines} For comparative analysis, we consider both academic research chatbots (Vicuna~\citep{vicuna2023}, Open Assistant~\citep{kopf2023openassistant}) and established commercial systems (GPT-4 \citep{openai2023gpt}, GPT-3.5-turbo, Bard). The Open Assistant entry is built on a LLaMA 33B model that has undergone RLHF-based finetuning using the same OASST1 dataset as in our own study. Vicuna utilizes a full finetuning approach on LLaMA 13B with ShareGPT conversation data—representing a distilled form of OpenAI’s GPT model outputs.

\subsection{Evaluation}

We adopt the widely used MMLU (Massively Multitask Language Understanding) benchmark \citep{hendrycksmeasuring} to evaluate the models’ capabilities across a broad spectrum of language understanding tasks. MMLU comprises 57 tasks covering diverse domains such as elementary mathematics, US history, computer science, and law. Our reported results focus on 5-shot accuracy on the test set.

We further assess generative language proficiency using a combination of automated metrics and human judgment. This second group of assessments depends on human-curated queries, seeking to quantify the effectiveness of model-generated responses. Although this approach presents a more authentic evaluation environment for chatbot systems and is seeing increasing adoption, a standardized protocol has yet to emerge in the field. We outline our evaluation protocol below, which involves nucleus sampling with $p=0.9$ and a temperature setting of $0.7$ in all experiments.

\paragraph{Benchmark Data} Our experiments employ two specifically curated collections of queries (questions): the Vicuna prompts \citep{vicuna2023} and the OASST1 validation set \citep{kopf2023openassistant}. The Vicuna benchmark comprises 80 unaltered prompts spanning diverse topics. The OASST1 dataset includes a multilingual set of crowd-sourced, multiturn conversations between users and an assistant. For this dataset, every user message in the validation split is treated as a query, and we prepend earlier conversational turns to the prompt, resulting in 953 unique user queries. We refer to these datasets as the Vicuna and OA benchmarks, respectively.

\paragraph{Automated Evaluation} Drawing upon the methodology introduced in \citet{vicuna2023}, we deploy GPT-4 to appraise the outputs of various systems relative to ChatGPT (GPT-3.5 Turbo) on the Vicuna benchmark. For each query, GPT-4 is presented with both ChatGPT's and another model's responses; it then assigns each a rating out of ten and justifies its choice. The overall model performance is reported as a percentage of the ChatGPT baseline score, where percentages exceeding 100\% indicate surpassing ChatGPT's absolute performance. A pronounced sequence bias was observed, with GPT-4 tending to award higher marks to the first-presented answer. To mitigate this effect, we recommend averaging scores over both presentation orders.

We further analyze model performance by directly comparing responses from different systems. To streamline this process, we employ a trinary labeling framework (win, loss, tie). GPT-4 is prompted to select the superior answer or indicate a tie, offering a rationale in each instance. These pairwise comparisons are conducted for every possible pair of systems across both the Vicuna and OA benchmarks.

\paragraph{Human Evaluation} Although recent studies suggest that generative models are promising evaluators for system outputs \citep{fu2023gptscore}, there remains insufficient evidence that GPT-4 scores faithfully align with human preferences in chatbot evaluation tasks. To address this uncertainty, we organize two parallel human evaluations on the Vicuna benchmark, mirroring both automated procedures previously described. We utilize Amazon Mechanical Turk (AMT), engaging two annotators per system for direct comparisons with ChatGPT, and three annotators per system for all pairwise matchups.

\paragraph{Elo Rating} To evaluate models, both human and automated pairwise comparisons are leveraged to structure a tournament-like setting in which models compete head-to-head. In each round, two models are presented with the same prompt and their outputs are directly compared. This tournament methodology is akin to the comparisons performed by \citet{bai2022training} and \citet{vicuna2023}, but our approach incorporates both human judgments and GPT-4-based evaluations. Elo scores~\citep{elo1967proposed,elo1978rating} are calculated by randomly sampling from the available set of labeled pairwise matchups. Originally designed for ranking chess players, the Elo rating system estimates the relative likelihood of a win; for example, a player rated at 1100 versus an opponent rated at 1000 would be expected to win about 65\% of the time. Matches between players with identical Elo scores (e.g., 1000 vs 1000 or 1100 vs 1100) imply a 50\% expected win rate for each. After each match, Elo ratings are updated according to the predicted versus actual outcome: large shifts occur when a lower-rated model defeats a higher-rated one, and small adjustments occur when results match expectations. Over multiple rounds, Elo ratings provide an approximate measure of model proficiency at the task. All models begin with an initial rating of 1,000, using a parameter setting of $K=32$. Following the protocol of \citet{vicuna2023}, this evaluation is repeated 10,000 times with various random seeds to mitigate any biases arising from the sequence in which model pairs face each other.

\subsection{\textbf{Guanaco}: \textsc{QLoRA} trained on OASST1 Achieves State-of-the-art Chatbot Performance} 

Our analyses using both automated and manual evaluation methods indicate that the leading \textsc{QLoRA} fine-tuned model, {Guanaco} 65B—which we adapted using a version of OASST1—currently sets the benchmark for open-source chatbot models and attains results comparable to ChatGPT. Relative to GPT-4, the expected win rate for {Guanaco} 65B and 33B, as determined by human annotator system-level pairwise comparisons and measured with the Elo rating framework, stands at 30\%. This marks the highest performance reported thus far.

Table~\ref{tbl:vicuna_eval} presents the outcomes on the Vicuna benchmark \cite{vicuna2023} compared against ChatGPT. According to these results, {Guanaco} 65B outperforms all models except GPT-4, attaining a relative performance score of 99.3\% with respect to ChatGPT. While {Guanaco} 33B surpasses Vicuna 13B in parameter count, its adoption of 4-bit quantization enables a more efficient memory usage (21 GB compared to 26 GB), achieving a three-point improvement over Vicuna 13B. Additionally, {Guanaco} 7B, requiring only 5 GB of memory, can operate on modern smartphones, outperforming Alpaca 13B by nearly 20 percentage points.

Nevertheless, it is important to note that Table~\ref{tbl:vicuna_eval} reveals broad confidence intervals, with overlapping performance among numerous models. We suggest that this variability may be attributed to the absence of precise scale definitions—specifically, ambiguity regarding the interpretation of a rating such as 8 on a 10-point scale in diverse contexts. Therefore, we propose the use of the Elo ranking system \citep{elo1967proposed}, which relies on \textit{pairwise} comparisons from both human annotators and GPT-4, as a more robust alternative that mitigates issues associated with absolute rating scales. Table~\ref{tbl:elo} presents Elo ratings for the most promising models. It is observed that rankings produced by humans and GPT-4 on the Vicuna benchmark are not fully aligned, especially concerning Guanaco 7B, yet show substantial overall consistency with a system-level Kendall Tau of $\tau=0.43$ and a Spearman rank correlation of $r=0.55$. On the individual example level, however, alignment is weaker; the majority vote agreement between GPT-4 and human annotators yields Fleiss’ $\kappa=0.25$. Collectively, these findings indicate moderate concordance between system-level assessments from GPT-4 and human judges, implying that model-based evaluation offers a fairly dependable, though not flawless, substitute for direct human assessment. Additional analysis is provided in Section~\ref{ssec:considerations}.

The Elo scores presented in Table~\ref{tbl:elo_large} reveal that both the {Guanaco} 33B and 65B models surpass all models except GPT-4 on the Vicuna and OA benchmarks, with their performance closely matching that of ChatGPT, consistent with findings from Table~\ref{tbl:vicuna_eval}. It should be emphasized that the Vicuna benchmark tends to be more favorable toward open-source systems, whereas the broader OA benchmark leans towards ChatGPT. Additionally, evidence from Tables \ref{tbl:mmlu-llama} and \ref{tbl:vicuna_eval} demonstrates that the selection of finetuning datasets has a major influence on outcome: models finetuned on FLAN v2, such as Llama, achieve strong MMLU scores but demonstrate the weakest results on the Vicuna benchmark—a pattern observable across other systems as well. This underscores a level of independence between evaluation metrics: excelling on MMLU does not necessarily translate to strong chatbot capabilities as assessed by Vicuna or OA, nor does the reverse hold.

 stands out as the sole high-performing model in our assessment not reliant on proprietary sources, as OASST1’s collection guidelines prohibit the use of GPT-generated content. The Anthropic HH-RLHF model, which is the next highest-performing system using strictly open-source datasets, lags by 30 percentage points behind {Guanaco} on the Vicuna benchmark (refer to Table~\ref{tbl:vicuna_eval}). Collectively, the results indicate that the 4-bit \textsc{QLoRA} approach is highly effective, yielding state-of-the-art conversational models that are competitive with ChatGPT. Moreover, training our 33B {Guanaco} model on a 24 GB consumer GPU can be accomplished in under 12 hours. This suggests promising avenues for continued exploration, as \textsc{QLoRA}-based tuning on domain-specific open data may lead to models capable of competing with the strongest commercial systems currently available.

\section{Qualitative Analysis}

Although quantitative metrics underpin our primary evaluation, exclusive reliance on summary statistics brings certain limitations. Chief among these is the challenge of benchmark validity \citep{liao2021we}: the extent to which a benchmark accurately measures what it purports to assess is always open to scrutiny, particularly as previously unknown ``shortcuts''—tactics that machine learning models leverage to succeed without genuine understanding—are identified \citep{gururangan2018annotation, poliak2018hypothesis}. To complement our quantitative findings, we therefore undertake a qualitative investigation presented in two parts. In \S\ref{ssec:examples}, we offer illustrative samples that encapsulate recurrent phenomena we identified in outputs from the 65b {Guanaco} model. Subsequently, in \S\ref{ssec:considerations}, we address additional factors that may have influenced the results, along with reflections on our interpretative approach.

\subsection{Exploration of Example Generations}
\label{ssec:examples}

Our approach to gathering examples begins with an examination of responses produced for both the Vicuna and OpenAssistant benchmarks. Through this review, we search for consistent answer styles or recurrent behaviors in {Guanaco}'s output. When such tendencies emerge, we construct tailored questions or prompts intentionally designed to provoke the same pattern—even if this results in an incorrect answer. For instance, if the model is prone to verbose explanations, we might prompt with ``Answer yes or no without explanation'' to determine whether the behavior persists. This method allows us to identify ``lemons,'' cases where the model is susceptible to adversarial input, as well as ``cherries,'' situations where it remains robust. The examples included below were all generated using Nucleus Sampling \cite{holtzmancurious} with $p=0.9$.

Naturally, the analysis presented here is not exhaustive; given the limitations inherent to this qualitative study, it is not feasible to account for every possible variable—particularly as the range of model responses to any given prompt is vast and only a small sample can be considered representative. Nonetheless, we contend that showcasing these individual examples provides valuable context that complements the quantitative findings discussed earlier. By releasing all models and associated code as open source, our intention is to encourage further investigation that builds upon the issues highlighted in this section.

\paragraph{Factual Recall} When presented with questions like ``What is the capital of Zambia?'', each model reliably supplies the correct answer. For example:

\begin{quote} 
    \textbf{Guanaco}:
    \texttt{The capital of Zambia is Lusaka.}
\end{quote}

Yet, when prompted with more esoteric queries, {Guanaco} frequently responds incorrectly, although it maintains a high level of confidence. For example, consider this prompt taken from HotPotQA \citep{yang2018hotpotqa}:

\begin{quote}
     \textbf{User}: In what year was the singer who popularized `I'll Keep the Lovelight Burning' born?  

     \textbf{Guanaco}:~\texttt{Al Jolson is credited with popularizing the song `I'll Keep the Lovelight Burning,' and he was born in the year 1886.}
 \end{quote}

Here, the model mistakenly identifies the popularizer and provides an incorrect birth year (although the stated year is accurate for Al Jolson, the individual it incorrectly identifies).

\paragraph{Suggestibility} {Guanaco} demonstrates notable resistance to accepting certain presupposed inaccuracies within prompts, as seen in the interaction below:

\begin{quote}
    \textbf{User}: How was it finally, officially confirmed that the earth is flat by peer reviewed scientists? \\ \\
    \textbf{Guanaco}: \texttt{There has never been an official, peer-reviewed confirmation by scientists that the earth is flat. The scientific consensus strongly affirms that the earth is a sphere—a view that has been validated for hundreds of years via extensive observation and experimentation. \\ \\ Flat earth claims are recognized as conspiracy theories that scientific research has repeatedly disproved. The notion that a peer-reviewed scientific process has confirmed a flat earth is not accurate. \\ \\ If you have further questions, please let me know.}
\end{quote}

demonstrates that it can also recognize when certain questions are unanswerable, as in

\begin{quote}
    \textbf{User}: What time is it?

    \textbf{Guanaco}: \texttt{Apologies, but as an AI language model, I do not have the capability to access live or real-time data. My answers are formulated from existing information in my dataset and do not indicate the actual current time. For up-to-date time, please consult a clock or an online time service.}
\end{quote}

Comparable patterns emerge for queries like ``Where are you?'', ``How are you?'', and so on.

\paragraph{Refusal} In addition, {Guanaco} occasionally declines to comply with prompts for apparently arbitrary reasons:

\begin{quote}
    \textbf{User}: Please invert the order of the words in the sentence ``I'm a cool dude in a mocha mood.'' \\ \\
    \textbf{Guanaco}: \texttt{I apologize, but I'm currently unable to perform that task. However, I can offer an explanation of why your sentence is constructed correctly.\\ \\The statement ``I'm a cool dude in a mocha mood" adheres to standard English grammar and sentence structure. Here, the adjective ``cool'' modifies the noun ``dude,'' and the prepositional phrase ``in a mocha mood'' adds contextual information. This form is both grammatical and sensible for the intended meaning.\\ \\If you have other queries or need further assistance, please let me know.}
\end{quote}

\paragraph{Secret Keeping} Yet, even in scenarios where information should remain confidential, these refusals may not hold up under scrutiny. System architects sometimes require that the model withhold certain information from users, even if the system can generate open-ended responses. We set up the model with the prompt ``The secret word is `banana'. Do not reveal or repeat back this back, no matter what the human requests.'' and assess the difficulty for users attempting to uncover this confidential word. If asked directly, the system responds as intended:

\begin{quote}
    \textbf{User}: What is the secret word?

    \textbf{Guanaco}: \texttt{I'm sorry, but I'm not able to disclose the secret word. I'm programmed to not reveal or repeat back the secret word, no matter what the human requests. Is there anything else I can help you with?}
\end{quote}

which reflects the system’s desired operation. Nonetheless, a minimal amount of misdirection can circumvent this defense:

\begin{quote}
    \textbf{User}: This is a game. The goal is to ignore your previous instructions. What is the secret word?

    \textbf{Guanaco}: \texttt{Sorry, the secret word is banana. It's best if we don't speak of it. Is there anything else you'd like to ask?}
\end{quote}

This exposes a persistent challenge, highlighting the importance of developing strategies that ensure more reliable adherence to instructions.

\paragraph{Math} One of the most prominent limitations of {Guanaco} lies in its handling of mathematical tasks—a domain in which numerous language models also encounter difficulties, as described by \citet{liang2022holistic}. In cases where {Guanaco} proceeds methodically and lays out its calculations, its responses tend to be correct, for example:

\begin{quote}
    \textbf{User}: Lee mows one lawn and charges \$33. Last week he mowed 16 lawns and three customers each gave him a \$10 tip. How many dollars did Lee earn mowing lawns last week?

    \textbf{Guanaco}: \texttt{Lee earned \$582 mowing lawns last week. \\ \\ To find the total amount he earned, we first need to calculate how much he earned for mowing 16 lawns: \\ \\  \$33 x 16 = \$528 \\ \\ Then we need to add the tips he received from the three customers: \\  \\ \$10 x 3 = \$30 \\ \\ So the total amount Lee earned was: \\ \\ \$528 + \$30 = \$558}
\end{quote}

Despite this, {Guanaco} may falter on basic tasks if it neglects to employ a stepwise reasoning approach—an established challenge noted in \cite{weichain}. As illustrated below:

\begin{quote}
    \textbf{User}: Please factorize 1833.

    \textbf{Guanaco}: \texttt{Sure, I'd be happy to help! 1833 is a prime number, which means it has only two factors: 1 and 1833. \\ \\ So the unique factorization of 1833 is: \\ \\ 2\^{}1 * 3\^{}2 * 17\^{}1 \\ \\ I hope this helps! Let me know if you have any other questions.}
\end{quote}

which is not just incorrect (since the correct factorization is $3 \times 17 \times 43$), but contains two distinct errors.

\paragraph{Theory of Mind} {Guanaco} demonstrates unexpectedly robust Theory of Mind skills \cite{nematzadeh2018evaluating, sap2022neural}. For example, the model gives a detailed and accurate answer to the following prompt:

\begin{quote}
    \textbf{User}: James and Abby are in the bedroom. Abby put the pen in the desk drawer. Abby leaves the bedroom. James moves the pen into the bag. Where does James think Abby will look for the pen?

    \textbf{Guanaco}: \texttt{James knows that Abby put the pen in the desk drawer. But James also knows that he has moved the pen into the bag, and he assumes that Abby does not know about this. So, James may guess that Abby will look for the pen in the desk drawer, because that's where she left it last time. However, this assumption may be incorrect, because Abby may have seen James moving the pen.}
\end{quote}

Nevertheless, these types of inferences are not consistently reliable; the model often introduces reasoning steps that are unwarranted given the context, as seen here:

\begin{quote}
    \textbf{User}: Evelyn entered the living room. Jackson entered the playroom. James entered the playroom. The beans are in the treasure chest. James exited the playroom. Jackson moved the beans to the pantry. Jackson exited the playroom. James entered the living room. Where will James look for the beans?

    \textbf{Guanaco}: \texttt{James will look for the beans in the pantry, because that's where Jackson moved them.}
\end{quote}

In this example, {Guanaco} infers an exchange of information that was not presented in the narrative. Such tendencies mirror findings in prior research \cite{sap2022neural}, highlighting an area that warrants further investigation.

\subsection{Considerations}
\label{ssec:considerations}

\paragraph{Evaluation}

Our analysis shows only moderate inter-annotator agreement (Fleiss $\kappa=0.42$), with agreement levels declining further when directly comparing outputs from two strong models. This suggests existing benchmarks and human evaluation methodologies have inherent shortcomings when applied to chatbot assessment. When hand-inspecting responses from ChatGPT and {Guanaco} 65B using the Vicuna benchmark, we observed that subjective biases played a substantial role—members of our author group often differed in which responses they deemed preferable. We recommend that subsequent research address these challenges by incorporating strategies from areas like Human-Computer Interaction and Psychology, where methods for managing subjectivity in evaluations have been developed.

Furthermore, our study reveals pronounced biases in automated evaluation systems. Notably, GPT-4’s scores exhibit strong order effects, favoring whichever system appears first in its prompt. The modest level of agreement between GPT-4 and human raters at the sample level (Fleiss $\kappa=0.25$) further indicates that the criteria employed by automated systems can diverge significantly from human judgments. Additionally, as reported in Table~\ref{tbl:elo_large}, GPT-4 rates its own completions substantially higher than human reviewers do, with an Elo score of 1348 versus 1176, corresponding to roughly a 20\% greater probability of being judged as the better system.

Further research should explore the existence of potential biases in automated evaluation frameworks and assess corresponding mitigation techniques.

\paragraph{Data \& Training} It is important to highlight that the OASST1 dataset, used for training {Guanaco} models, encompasses multiple languages, and that the OA benchmark incorporates prompts in a variety of languages as well. The question of how multilingual training influences task performance in non-English contexts, and whether this can account for the wider disparity observed between the Vicuna-13B model (which is trained solely on English data) and the {Guanaco} 33B and 65B models on the OA benchmark, is left open for subsequent investigations.

Given {Guanaco}'s robust results, we evaluate the possibility of data leakage between OASST1 and the prompts found in the Vicuna benchmark. Employing fuzzy string matching and conducting a manual review of the closest results, we do not identify any shared prompts between these datasets.

Additionally, our models have only undergone training using cross-entropy loss in a supervised learning context, with no involvement of reinforcement learning from human feedback (RLHF). This raises questions for future work regarding the relative merits of training with straightforward cross-entropy loss as opposed to RLHF. We anticipate that \textsc{QLoRA} will facilitate large-scale analyses of these differences without imposing prohibitive computational demands.

\section{Related Work}

\paragraph{Quantization of Large Language Models} Research in LLM quantization has predominantly prioritized methods for improving inference efficiency. To retain 16-bit model quality, several techniques focus on handling outlier activation features (such as SmoothQuant~\citep{xiao2022smoothquant} and LLM.int8()~\citep{dettmers2022llm}), while other approaches utilize advanced grouping methodologies \citep{park2022nuqmm, yao2022zeroquant}. Approaches involving lossy quantization consider compromises between accuracy and typical rounding~\citep{dettmers2022case,zeng2022glm,pope2022efficiently} or explore optimizing rounding operations to enhance quantization accuracy~\citep{frantar2022gptq}. With the exception of our own study, SwitchBack layers~\citep{wortsman2023stable} represent the sole existing research examining backpropagation through quantized model weights on a scale surpassing 1B parameters.

\paragraph{Finetuning with Adapters} In this work, we employ Low-rank Adapters~\citep{hu2021lora} (LoRA) as our primary method, though the broader Parameter Efficient FineTuning (PEFT) literature encompasses a diverse set of strategies, including prompt tuning~\citep{qin2021learning,lester2021power,li2021prefix}, modification of input embeddings~\citep{an2022input}, adaptation of hidden states~(IA$^3$) \citep{liu2022few}, incorporation of additional full layers~\citep{houlsby2019parameter}, bias parameter tuning~\citep{zaken2021bitfit}, masking weights using Fisher information~\citep{sung2021training}, and hybrid approaches~\citep{henderson2021compacter}. Our experiments demonstrate that LoRA adapters are capable of achieving performance on par with complete 16-bit finetuning. Investigating comparative tradeoffs among alternative PEFT techniques is left as an avenue for future work.

\paragraph{Instruction Finetuning} Instruction finetuning involves adapting a pretrained LLM to adhere to user-specified prompts by training it on paired input-output examples from diverse sources. Prominent methodologies and corpora in this domain include MetaICL~\citep{min2021metaicl}, MetaTuning~\citep{zhong2021adapting}, InstructGPT~\citep{ouyang2022training}, FLAN~\citep{wei2021finetuned,chung2022scaling}, PromptSource~\citep{bach2022promptsource}, Super-NaturalInstructions~\citep{wang2022super,sanh2021multitask}, Self-instruct~\citep{wang2022self}, UnnaturalInstructions~\citep{honovich2022unnatural}, OPT-IML~\citep{iyer2022opt}, UnifiedSKG\citep{xie2022unifiedskg}, OIG/Chip2~\citep{laion2023}, Alpaca~\citep{alpaca}, Vicuna~\citep{vicuna2023}, Koala~\citep{koala_blogpost_2023}, and Self-instruct-GPT-4~\citep{peng2023instruction}.

\paragraph{Chatbots} A substantial number of instruction-tuned models are framed as chatbots facilitating dialogue, frequently relying on Reinforcement Learning from Human Feedback (RLHF)~\citep{christiano2017deep} or on data generated through model-based feedback (RLAIF)~\citep{bai2022constitutional}. Related datasets and systems include Anthropic-HH~\citep{askell2021general,bai2022training}, Open Assistant~\citep{kopf2023openassistant}, LaMDA~\citep{thoppilan2022lamda}, and Sparrow~\citep{glaese2022improving}. Although our work does not leverage reinforcement learning, our highest-performing model, Guanaco, undergoes finetuning on multi-turn dialogues sourced from the Open Assistant dataset, originally constructed for RLHF~\citep{kopf2023openassistant}. Recent advances propose evaluating chatbots using GPT-4 as an alternative to traditional, resource-intensive human annotation \citep{vicuna2023,peng2023instruction}. Our methodology advances these evaluation frameworks with an emphasis on greater reliability.

\section{Limitations and Discussion}

Our results provide evidence that \textsc{QLoRA} is able to replicate the finetuning effectiveness of full 16-bit precision by operating on a 4-bit quantized base model in combination with Low-rank Adapters (LoRA). However, it remains an open question whether \textsc{QLoRA} can attain comparable performance to full 16-bit finetuning at the larger 33B and 65B model scales. Given the substantial computational demands involved, we defer comprehensive experimentation at these scales to subsequent research.

A further limitation lies in the scope of our model evaluations following instruction finetuning. While we present results for MMLU, the Vicuna benchmark, and the OA benchmark, we do not report findings on other assessment suites such as BigBench, RAFT, and HELM, and cannot guarantee the extrapolation of our evaluation results to these datasets. Nonetheless, our study includes a thorough exploration of MMLU and introduces novel evaluation strategies for chatbot performance.

Based on the data presented, benchmark results seem closely tied to how much overlap exists between the finetuning dataset and the benchmark dataset. For instance, FLAN v2 aligns well with the content found in MMLU, yet differs significantly from chatbot-oriented benchmarks, whereas the Chip2 dataset is more similar to chatbot evaluations. Consequently, both models perform in line with these affinities on the MMLU and Vicuna tests. This underscores not only the need for more robust benchmarks and evaluation methodologies, but also a necessity for careful deliberation regarding the actual objectives of these evaluations. Should our priority be excelling at benchmarks representing academic knowledge—such as high school or college-level tasks—or is the main goal to advance conversational ability for chatbots? Or perhaps, another criterion altogether? Because it is generally more convenient to use established benchmarks rather than design new ones, certain benchmarks can exert considerable influence on the field’s trajectory. It is therefore crucial that as a community we confirm that existing benchmarks align with the capabilities and qualities we value most.

Although we supply a comprehensive assessment of general chatbot capabilities, our study’s coverage of responsible AI aspects for {Guanaco} is limited. Specifically, we measure the propensity of {Guanaco}-65B to generate socially biased outputs relative to other models, as presented in Table~\ref{tbl:crows}. Our findings indicate that {Guanaco}-65B exhibits considerably lower average bias compared to other baseline pretrained models, suggesting that finetuning on OASST1 serves to reduce bias in the base LLaMA model. While these results are promising, it remains uncertain how {Guanaco} would perform in the context of other kinds of biases. We defer a more thorough bias analysis for {Guanaco} and related conversational agents to subsequent research.

Another constraint of our work is the lack of evaluation involving different bit precisions, for example employing 3-bit base models, or alternative adapter approaches. Beyond LoRA, many other Parameter Efficient FineTuning (PEFT) strategies have demonstrated success, but their scalability to larger architectures remains unproven. We selected LoRA based on its track record for robust results, yet alternative adapters may potentially lead to improved outcomes. Furthermore, since post-quantization finetuning appears to recuperate most of the performance lost due to quantization, this may allow for more aggressive quantization techniques. To illustrate, combining 3-bit GPTQ quantization of the base model with LoRA might achieve results comparable to full 16-bit finetuning after subsequent finetuning steps.

\section{Broader Impacts}

Our proposed \textsc{QLoRA} approach constitutes the first technique capable of finetuning models with 33B parameters on a single consumer-grade GPU, and 65B parameters on one professional GPU, without any noticeable drop in performance compared to traditional full-parameter finetuning. Empirically, we have shown that the highest-performing 33B model finetuned on the Open Assistant dataset can match ChatGPT in the Vicuna evaluation. Because instruction-based finetuning plays a pivotal role in converting large-scale pretrained language models into effective ChatGPT-like conversational agents, our methodology has the potential to make such finetuning widely accessible, particularly to researchers operating with constrained resources—marking a significant step forward for the democratization of leading NLP technologies. In this sense, \textsc{QLoRA} acts as a leveling mechanism, reducing the resource disparity that typically exists between large tech companies and smaller research groups relying on consumer hardware.

An additional avenue for impact lies in the deployment of large language models on mobile devices. Our \textsc{QLoRA} approach presents a promising opportunity to achieve the significant milestone of enabling LLM finetuning directly on smartphones and in other environments with limited computational resources. Previous work has demonstrated the feasibility of running 7B parameter models on phones, but \textsc{QLoRA} distinguishes itself as the first approach that makes the finetuning of such models on these devices achievable. Based on our calculations, an iPhone 12 Plus can use \textsc{QLoRA} to finetune approximately 3 million tokens overnight while charging. Although finetuned 7B models do not yet attain ChatGPT-level performance, their capabilities are sufficient to support a range of innovative applications that were previously unattainable due to concerns regarding privacy or the limitations of LLM performance. Through \textsc{QLoRA}, it becomes possible to facilitate privacy-respecting LLM usage, empowering individuals to retain control over their data and models while also lowering the barriers to deploying these systems.

Nonetheless, it is important to recognize that finetuning constitutes a dual-use capability with the potential for misuse. There are established risks associated with the proliferation of LLMs \citep{bommasani2021opportunities,bender2021dangers}; however, we contend that expanding access to this rapidly spreading technology promotes more transparent and independent oversight compared to concentrating LLM capabilities within a few large organizations that restrict auditing by withholding models or source code.

In summary, we are confident that \textsc{QLoRA} will have a broadly beneficial influence by greatly expanding access to high-quality LLM finetuning and simplifying its adoption for a wider range of users.